\chapter*{Lời mở đầu}
\addcontentsline{toc}{chapter}{Lời mở đầu}

Trong chương trình Hóa học đại cương, đặc biệt trong môi trường đào tạo kỹ thuật và nghiên cứu của Trường Đại học Bách khoa, việc kết hợp giữa lý thuyết và thực nghiệm là yếu tố then chốt để hiểu rõ bản chất của các quá trình biến đổi vật chất. Không chỉ dừng lại ở việc tiếp nhận kiến thức trên giảng đường, sinh viên còn cần trực tiếp quan sát, đo đạc và xử lý số liệu để hình thành tư duy khoa học, khả năng phân tích và tác phong làm việc cẩn trọng trong phòng thí nghiệm.

Báo cáo này trình bày kết quả thực hành của ba nội dung trọng tâm trong học phần, gồm nhiệt động học, động hóa học và hóa học phân tích. Thông qua Bài 2 (Nhiệt phản ứng), chúng em khảo sát sự trao đổi năng lượng và định luật bảo toàn nhiệt năng trong các phản ứng hóa học, từ đó làm rõ bản chất của hiệu ứng nhiệt. Tiếp đó, Bài 4 (Bậc phản ứng) giúp phân tích cơ chế và tốc độ phản ứng bằng cách xác định sự phụ thuộc của thời gian phản ứng vào nồng độ các chất tham gia. Cuối cùng, Bài 8 (Phân tích thể tích) rèn luyện kỹ năng định lượng chính xác thông qua phương pháp chuẩn độ axit--bazơ, đồng thời củng cố cách xử lý kết quả thực nghiệm theo hướng định lượng.

Chuỗi thí nghiệm này không chỉ góp phần củng cố kiến thức nền tảng mà còn giúp sinh viên Bách khoa rèn luyện tư duy phân tích, kỹ năng thao tác chuẩn xác và tinh thần làm việc nghiêm túc trong nghiên cứu thực nghiệm. Thông qua quá trình thực hiện, đối chiếu lý thuyết với số liệu thu được và đánh giá sai số, báo cáo hướng đến việc phản ánh rõ hơn vai trò của thực hành trong việc xây dựng nền tảng chuyên môn vững chắc cho các chuyên ngành sau này ở bậc đại học.